%%%%%%%%%%  scan idx 221  |  folio 208  %%%%%%%%%%
\emph{Preuve.} D'après (T.F. II 4.17.1), on a
\[
  R^i p_*(E)_s=H^i\bigl(X_s,i_s^{-1}(E)\bigr).
\]
Mais, comme $i_s^{-1}(\cO_X)$ est mou (2.3), on a
\[
  H^i\bigl(X_s,i_s^{-1}(E)\bigr)=0\quad \text{pour } i>0,
\]
d'où l'assertion.

\textbf{Corollaire 2.6.} Soient $X,Y$ des $S$-variétés mixtes, et $s\in S$. On suppose que $X$ satisfait aux hypothèses de (2.5). Pour tout morphisme $f:X_s\to Y_s$, il existe un voisinage ouvert $S'$ de $s$ et un morphisme $\bar f:X|S'\longrightarrow Y|S'$ tel que $\bar f_s=f$.

\emph{Preuve.} Conséquence immédiate de (2.3) (avec $A=X_s$) et du fait que, l'on a, grâce aux hypothèses sur $X$,
\[
  H^0\bigl(X_s,i_s^{-1}\underline{\Homop}(X,Y)\bigr)
  =\operatorname*{Lim}_{S'\ni s}\Homop(X|S',Y)\qquad (S'\ \text{ouvert})
\]
(T.F. II 3.3.1).

\textbf{Proposition 2.7.} Soient $X,Y$ des $S$-variétés mixtes, $f:X\to Y$ un morphisme, $s\in S$, $x\in X_s$. Si l'application linéaire tangente à $f_s:X_s\to Y_s$ est inversible en $x$, alors $f$ est un isomorphisme local en $x$, i.e. il existe un voisinage ouvert $U$ de $x$ dans $X$ et un ouvert $V$ de $Y$ tels que $f$ induise un isomorphisme de $U$ sur $V$.

\emph{Preuve.} C'est une conséquence facile du théorème des fonctions implicites. Les détails sont laissés au lecteur.

\textbf{Corollaire 2.8.} Soient $X,Y$ des $S$-variétés mixtes, et $s\in S$. On suppose que $s$ admet un système fondamental de voisinages $S'$ tels que les $X|S'$ (resp. $Y|S'$) forment un système fondamental de voisinages de $X_s$ (resp. $Y_s$) (condition réalisée par exemple si $S$ est localement compact et $X$ (resp. $Y$) propre sur $S$). Soit $f:X\to Y$ un $S$-morphisme tel que $f_s:X_s\to Y_s$ soit un isomorphisme. Il existe alors un voisinage $S'$ 
